\section*{अध्याय \ref{ch:b2:measurement-statistics} --- रासायनिक मापन की सांख्यिकी}

\begin{solution}{exo:b2:measurement-statistics:1}
$\bar V = 62.41/5 = \qty{12.482}{mL}$; $\sum(V_i - \bar V)^2 = 0.00328$, $s = \sqrt{0.00328/4} \approx
\qty{0.029}{mL}$; $s/\sqrt5 \approx \qty{0.013}{mL}$।
\end{solution}

\begin{solution}{exo:b2:measurement-statistics:2}
4 स्वातंत्र्य कोटियों के लिए $t = 2.776$: $12.482 \pm 2.776 \times 0.0128$, अर्थात् $(12.48 \pm 0.04)$~mL।
\end{solution}

\begin{solution}{exo:b2:measurement-statistics:3}
द्रव्यमान और आयतन की आपेक्षिक अनिश्चितताएँ 0.04~\% और 0.06~\% हैं; वर्ग-योग में
संयोजित करने पर,
\[\frac{u(c)}{c} = \sqrt{\Bigl(\frac{0.0002}{0.5000}\Bigr)^2 + \Bigl(\frac{0.15}{250.0}\Bigr)^2}
\approx 0.07~\%.\]
\end{solution}

\begin{solution}{exo:b2:measurement-statistics:4}
$\bar x = 1.5$, $\bar y = 0.31$, $S_{xx} = 5$, $S_{xy} = 0.99$: $b = 0.198$, $a = 0.31 - 0.198 \times 1.5
= 0.013$।
\end{solution}

\begin{solution}{exo:b2:measurement-statistics:5}
$|10.12 - 10.00|/(0.08/\sqrt5) = 0.12/0.036 \approx 3.4 > 2.776$: अंतर सार्थक है, विधि
अभिनत है (लगभग \qty{+0.12}{mg/L} से)।
\end{solution}

\begin{solution}{exo:b2:measurement-statistics:6}
$\mathrm{pH} = -\log_{10}c$, $\mathrm d\,\mathrm{pH}/\mathrm dc = -1/(c\ln10)$:
$u(\mathrm{pH}) = (u(c)/c)/\ln10 = 0.02/2.303 \approx 0.009$।
\end{solution}

\begin{solution}{exo:b2:measurement-statistics:7}
$x_{\mathrm{LOD}} = 3 \times 0.0012/0.0098 \approx \qty{0.37}{\micro g/L}$; $x_{\mathrm{LOQ}} = 10
\times 0.0012/0.0098 \approx \qty{1.2}{\micro g/L}$।
\end{solution}

\begin{solution}{exo:b2:measurement-statistics:8}
रेखा मापे गए विलयनों में \qty{3.11}{mg/L} देती है; नमूना पाँच गुना तनु किया गया था:
$3.11 \times 50.0/10.0 \approx \qty{15.6}{mg/L}$।
\end{solution}

\begin{solution}{exo:b2:measurement-statistics:9}
\qty{0.05}{mg/L} \omterm{def:b2:measurement-statistics:validation}{पुनरावर्तनीयता} है, \qty{0.15}{mg/L} \omterm{def:b2:measurement-statistics:validation}{पुनरुत्पादनीयता}। हर प्रयोगशाला की अपनी
छोटी \omterm{def:b2:measurement-statistics:validation}{अभिनति} होती है (अंशांकन, मानक, उपकरण); ये \omterm{def:b2:measurement-statistics:validation}{अभिनतियाँ} प्रयोगशालाओं के बीच भिन्न होती हैं और
प्रकीर्णन में जुड़ जाती हैं, अतः \omterm{def:b2:measurement-statistics:validation}{पुनरुत्पादनीयता} सदैव बड़ी होती है।
\end{solution}

\begin{solution}{exo:b2:measurement-statistics:10}
$\partial S/\partial a = 0$ का अर्थ है $\sum(y_i - a - bx_i) = 0$: \omterm{def:b2:measurement-statistics:calibration}{अवशिष्टों} का योग शून्य है।
$n$ से भाग देने पर: $\bar y = a + b\bar x$, अतः $(\bar x, \bar y)$ रेखा पर है।
\end{solution}

\begin{solution}{exo:b2:measurement-statistics:11}
एक चरण:
\[\sqrt{(0.006/1.000)^2 + (0.08/100.0)^2} \approx 0.61~\%.\]
एक दस-गुना चरण:
\[\sqrt{(0.02/10.00)^2 + (0.08/100.0)^2} \approx 0.22~\%,\]
और ऐसे दो चरण वर्ग-योग में $0.22\sqrt2 \approx 0.30~\%$ देते हैं। दो-चरणीय मार्ग दुगुना परिशुद्ध है: एक-चरणीय मार्ग में
छोटा पिपेट प्रभावी है।
\end{solution}

\begin{solution}{exo:b2:measurement-statistics:12}
$z = 1.6/\sqrt{0.4^2 + 0.5^2} \approx 2.5 > 2$: संगत नहीं; (कम-से-कम) एक प्रयोगशाला में \omterm{def:b2:measurement-statistics:validation}{अभिनति} है।
विस्तारित अनिश्चितताओं ($k = 2$) के साथ पहली $9.6 \pm 0.8$ देती है, जिसमें 10 समाविष्ट है; दूसरी
$11.2 \pm 1.0$, पूर्णतः 10 से ऊपर। जल पर कोई भी निर्णय देने से पहले असहमति को सुलझाना होगा,
उदाहरण के लिए किसी संदर्भ पदार्थ से।
% ledger: who:lead
\end{solution}

\begin{solution}{pb:b2:measurement-statistics:1}\emergencystretch=3em
\textbf{1.} अम्ल परमाणुकरण और अनुक्रिया को बदलता है; ढाल लागू होने के लिए मानकों और नमूनों की
आधात्री समान होनी चाहिए।
\textbf{2.} $\bar x = 10$, $\bar y = 0.1004$, $S_{xx} = 250$, $S_{xy} = 2.445$: $b = 0.00978$ प्रति
\unit{\micro g/L}, $a = 0.1004 - 0.0978 = 0.0026$।
\textbf{3.} $+0.0004$, $-0.0005$, $+0.0006$, $-0.0013$, $+0.0008$।
\textbf{4.} चिह्न एकांतर हैं, कोई वक्रता नहीं: 0--\qty{20}{\micro g/L} पर रेखा पर्याप्त है।
\textbf{5.} $s_{y/x} = \sqrt{3.1 \times 10^{-6}/3} \approx 0.00102$।
\textbf{6.} रिक्त (अम्ल, जल, भट्टी) की अपनी एक छोटी अवशोषकता होती है।
\textbf{7.} संवेदनशीलता: लेड के प्रति \unit{\micro g/L} अवशोषकता।
\textbf{8.} $\bar y_0 = 0.1010$; $s = 0.0030$।
\textbf{9.} $x_0 = (0.1010 - 0.0026)/0.00978 \approx \qty{10.06}{\micro g/L}$।
\textbf{10.} $s_{x_0} = (0.00102/0.00978)\sqrt{1/3 + 1/5 + (0.1010 - 0.1004)^2/(0.00978^2 \times 250)}
\approx 0.104 \times 0.730 \approx \qty{0.076}{\micro g/L}$।
\textbf{11.} $n - 2 = 3$; $t = 3.182$।
\textbf{12.} $10.06 \pm 3.182 \times 0.076 = (10.06 \pm 0.24)$~\unit{\micro g/L}।
\textbf{13.} पद $1/m$ 1 से घटकर $1/3$ हो जाता है: $s_{x_0}$ लगभग एक-तिहाई घट जाता है।
\textbf{14.} $f = 50.50/50.0 = 1.010$; $10.06 \times 1.010 \approx \qty{10.16}{\micro g/L}$।
\textbf{15.} $f = 1 + V_2/V_1$: $u^2(f) = (u(V_2)/V_1)^2 + (V_2u(V_1)/V_1^2)^2 = (0.005/50.0)^2 +
(0.50 \times 0.03/2500)^2$, $u(f) \approx 1.0 \times 10^{-4}$, आपेक्षिक रूप से 0.01~\%।
\textbf{16.} नहीं: 0.01~\%, जबकि $0.076/10.06 \approx 0.75~\%$।
\textbf{17.} $s_{\mathrm{blank}} \approx 0.00115$; $x_{\mathrm{LOD}} = 3 \times 0.00115/0.00978 \approx
\qty{0.35}{\micro g/L}$।
\textbf{18.} $10 \times 0.00115/0.00978 \approx \qty{1.2}{\micro g/L}$।
\textbf{19.} हाँ, बहुत अधिक।
\textbf{20.} हाँ: नल के जल में 9.92 से \qty{10.40}{\micro g/L} तक।
\textbf{21.} $(0.104 - 0.0026)/0.00978 \times 1.010 \approx \qty{10.5}{\micro g/L}$, और, अकेले पढ़ने पर,
एक आश्वस्त ``दिशानिर्देश से ऊपर'' जिसका आँकड़े समर्थन नहीं करते।
\textbf{22.} अंतराल बनाने में प्रयुक्त नियम उन श्रृंखलाओं में से 95~\% में वास्तविक मान को पकड़ता है जिन पर
उसे लागू किया जाता है; यह इस एक मान के बारे में कोई प्रायिकता नहीं है।
\textbf{23.} अधिक पुनरावृत्त पाठ्यांक और \qty{10}{\micro g/L} के निकट अधिक मानक (पद $1/m$ और
$1/n$), या अधिक संवेदनशील विधि; अंतराल केवल धीरे-धीरे सिकुड़ता है, $1/\sqrt{\text{संख्या}}$ के अनुसार।
\textbf{24.} किसी प्रमाणित संदर्भ जल का विश्लेषण करके, या नमूने में लेड की ज्ञात मात्रा मिलाकर
पुनर्प्राप्ति की जाँच करके।
\textbf{25.} \qty{10}{\micro g/L} के निकट सामान्य विधियों की अनिश्चितता कुछ प्रतिशत है, जैसा यहाँ है: कोई
निम्नतर दिशानिर्देश मान सामान्य प्रयोगशालाओं द्वारा विश्वसनीय रूप से जाँचा नहीं जा सकता।
\textbf{26.} 95~\% पर $\boldsymbol{(10.16 \pm 0.24)}$~\unit{\micro g/L}: अंतराल में
\qty{10}{\micro g/L} समाविष्ट है, अतः मापन न यह दिखाता है कि जल दिशानिर्देश से अधिक है और न यह कि वह
किसी अंतर से उसका पालन करता है।
\end{solution}
